% GENERATED FILE. Edit canonical lessons, then run npm run generate:books.
% canonical-source-hash: fnv1a64:b8b01202589dcce6
% canonical-lessons: UR-C344-barabar, UR-C344-imtiyaz, UR-C344-janibdari, UR-C344-zulm, UR-C344-uksana

\chapter{Fair and Unfair}
\label{ch:ur-fair-and-unfair}
\hlchaptermodality{344}

\begin{chapteropening}
\textbf{By the end of this chapter:} \emph{I can say whether people are treated fairly.}

Complete the last lesson of chapter 344: I can say whether people are treated fairly.
\end{chapteropening}

\section[\texorpdfstring{\ur{برابر} (barābar)}{برابر (barābar)}]{\ur{برابر} (barābar) — equal}

\label{lesson:UR-C344-barabar}



\begin{quote}
Before the new one: say the Urdu for an instruction, then the Urdu for to suggest.
\end{quote}

\subsection*{You'll want to know: \ur{برابر}}
\textbf{\ur{برابر}} — \emph{barābar} — \textquotedblleft{}equal\textquotedblright{}.

\begin{quote}
\ur{سب} \ur{انسان} \ur{برابر} \ur{ہیں۔}
\end{quote}

\begin{quote}
\emph{sab insān barābar haiṅ.}
\end{quote}

\begin{quote}
All human beings are equal.
\end{quote}

\begin{cousinweb}
\textbf{\ur{برابر}} (\emph{barābar}) is usually traced to Persian \emph{barābar}, \textquotedblleft{}equal, level\textquotedblright{}.
\end{cousinweb}

\begin{grammarlens}[title={What you've built}]
One more word for giving reasons and opinions.
\end{grammarlens}

\subsection*{Your turn}
\begin{itemize}
\raggedright
  \item \emph{Say it:} \emph{barābar}
  \item \emph{Say it:} \emph{barābar}, once more
  \item \emph{Say it:} \emph{sab insān barābar haiṅ.} — All human beings are equal.
  \item \emph{Recall:} say the Urdu for an instruction, then the Urdu for to suggest, then say \emph{barābar} again
\end{itemize}

\begin{tcolorbox}[breakable,colback=teal!4,colframe=teal!35!black,title={Before you move on}]
What does \ur{برابر} mean? (\textquotedblleft{}Equal\textquotedblright{}.) Say it once more.
\end{tcolorbox}
\section[\texorpdfstring{\ur{امتیاز} (imtiyāz)}{امتیاز (imtiyāz)}]{\ur{امتیاز} (imtiyāz) — discrimination; a distinction}

\label{lesson:UR-C344-imtiyaz}



\begin{quote}
Before the new one: say the Urdu for to suggest, then the Urdu for equal.
\end{quote}

\subsection*{You'll want to know: \ur{امتیاز}}
\textbf{\ur{امتیاز}} — \emph{imtiyāz} — \textquotedblleft{}discrimination; a distinction\textquotedblright{}.

\begin{quote}
\ur{ہمارے} \ur{اسکول} \ur{میں} \ur{کوئی} \ur{امتیاز} \ur{نہیں} \ur{ہوتا۔}
\end{quote}

\begin{quote}
\emph{hamāre iskūl meṅ koī imtiyāz nahīṅ hotā.}
\end{quote}

\begin{quote}
There is no discrimination in our school.
\end{quote}

Masculine.

\begin{cousinweb}
\textbf{\ur{امتیاز}} (\emph{imtiyāz}) is usually traced, through Persian, to Arabic \emph{imtiyāz}, \textquotedblleft{}being set apart, distinction\textquotedblright{}.
\end{cousinweb}

\begin{grammarlens}[title={What you've built}]
One more word for giving reasons and opinions.
\end{grammarlens}

\subsection*{Your turn}
\begin{itemize}
\raggedright
  \item \emph{Say it:} \emph{imtiyāz}
  \item \emph{Say it:} \emph{imtiyāz}, once more
  \item \emph{Say it:} \emph{hamāre iskūl meṅ koī imtiyāz nahīṅ hotā.} — There is no discrimination in our school.
  \item \emph{Recall:} say the Urdu for to suggest, then the Urdu for equal, then say \emph{imtiyāz} again
\end{itemize}

\begin{tcolorbox}[breakable,colback=teal!4,colframe=teal!35!black,title={Before you move on}]
What does \ur{امتیاز} mean? (\textquotedblleft{}Discrimination; a distinction\textquotedblright{}.) Say it once more.
\end{tcolorbox}
\section[\texorpdfstring{\ur{جانبداری} (jānibdārī)}{جانبداری (jānibdārī)}]{\ur{جانبداری} (jānibdārī) — bias; favouritism}

\label{lesson:UR-C344-janibdari}



\begin{quote}
Before the new one: say the Urdu for equal, then the Urdu for discrimination.
\end{quote}

\subsection*{You'll want to know: \ur{جانبداری}}
\textbf{\ur{جانبداری}} — \emph{jānibdārī} — \textquotedblleft{}bias; favouritism\textquotedblright{}.

\begin{quote}
\ur{اس} \ur{فیصلے} \ur{میں} \ur{جانبداری} \ur{تھی۔}
\end{quote}

\begin{quote}
\emph{is faisle meṅ jānibdārī thī.}
\end{quote}

\begin{quote}
There was bias in this decision.
\end{quote}

Feminine.

\begin{cousinweb}
\textbf{\ur{جانبداری}} (\emph{jānibdārī}) is usually traced to Persian \emph{jānibdārī}, \textquotedblleft{}taking a side\textquotedblright{}: Arabic \emph{jānib}, \textquotedblleft{}a side\textquotedblright{}, with Persian \emph{-dārī}, \textquotedblleft{}holding\textquotedblright{}.
\end{cousinweb}

\begin{grammarlens}[title={What you've built}]
One more word for giving reasons and opinions.
\end{grammarlens}

\subsection*{Your turn}
\begin{itemize}
\raggedright
  \item \emph{Say it:} \emph{jānibdārī}
  \item \emph{Say it:} \emph{jānibdārī}, once more
  \item \emph{Say it:} \emph{is faisle meṅ jānibdārī thī.} — There was bias in this decision.
  \item \emph{Recall:} say the Urdu for equal, then the Urdu for discrimination, then say \emph{jānibdārī} again
\end{itemize}

\begin{tcolorbox}[breakable,colback=teal!4,colframe=teal!35!black,title={Before you move on}]
What does \ur{جانبداری} mean? (\textquotedblleft{}Bias; favouritism\textquotedblright{}.) Say it once more.
\end{tcolorbox}
\section[\texorpdfstring{\ur{ظلم} (zulm)}{ظلم (zulm)}]{\ur{ظلم} (zulm) — oppression; cruelty}

\label{lesson:UR-C344-zulm}



\begin{quote}
Before the new one: say the Urdu for discrimination, then the Urdu for favouritism.
\end{quote}

\subsection*{You'll want to know: \ur{ظلم}}
\textbf{\ur{ظلم}} — \emph{zulm} — \textquotedblleft{}oppression; cruelty\textquotedblright{}.

\begin{quote}
\ur{کمزوروں} \ur{پر} \ur{ظلم} \ur{نہیں} \ur{ہونا} \ur{چاہیے۔}
\end{quote}

\begin{quote}
\emph{kamzoroṅ par zulm nahīṅ honā chāhiye.}
\end{quote}

\begin{quote}
The weak should not be oppressed.
\end{quote}

Masculine.

\begin{cousinweb}
\textbf{\ur{ظلم}} (\emph{zulm}) is usually traced, through Persian, to Arabic \emph{ẓulm}, \textquotedblleft{}wrongdoing, injustice\textquotedblright{}.
\end{cousinweb}

\begin{grammarlens}[title={What you've built}]
One more word for giving reasons and opinions.
\end{grammarlens}

\subsection*{Your turn}
\begin{itemize}
\raggedright
  \item \emph{Say it:} \emph{zulm}
  \item \emph{Say it:} \emph{zulm}, once more
  \item \emph{Say it:} \emph{kamzoroṅ par zulm nahīṅ honā chāhiye.} — The weak should not be oppressed.
  \item \emph{Recall:} say the Urdu for discrimination, then the Urdu for favouritism, then say \emph{zulm} again
\end{itemize}

\begin{tcolorbox}[breakable,colback=teal!4,colframe=teal!35!black,title={Before you move on}]
What does \ur{ظلم} mean? (\textquotedblleft{}Oppression; cruelty\textquotedblright{}.) Say it once more.
\end{tcolorbox}
\section[\texorpdfstring{\ur{اکسانا} (uksānā)}{اکسانا (uksānā)}]{\ur{اکسانا} (uksānā) — to provoke; to stir up}

\label{lesson:UR-C344-uksana}



\begin{quote}
Before the new one: say the Urdu for favouritism, then the Urdu for oppression.
\end{quote}

\subsection*{You'll want to know: \ur{اکسانا}}
\textbf{\ur{اکسانا}} — \emph{uksānā} — \textquotedblleft{}to provoke; to stir up\textquotedblright{}.

\begin{quote}
\ur{اس} \ur{نے} \ur{لوگوں} \ur{کو} \ur{ظلم} \ur{کے} \ur{خلاف} \ur{اکسایا۔}
\end{quote}

\begin{quote}
\emph{us ne logoṅ ko zulm ke khilāf uksāyā.}
\end{quote}

\begin{quote}
He stirred people up against oppression.
\end{quote}

Past, with \ur{نے} (he or she …): \ur{اس} \ur{نے} \ur{بھائی} \ur{کو} \ur{اکسایا}, \emph{us ne bhāī ko uksāyā}.

\begin{cousinweb}
\textbf{\ur{اکسانا}} (\emph{uksānā}) is inherited: CDIAL lists the Hindi cousin \emph{uksānā}, \textquotedblleft{}to draw out (a wick), to excite\textquotedblright{}, under Sanskrit \emph{útkasati}, \textquotedblleft{}opens\textquotedblright{} (CDIAL 1718). As a wick is drawn up to make a flame burn higher.
\end{cousinweb}

\begin{grammarlens}[title={What you've built}]
One more word for giving reasons and opinions.
\end{grammarlens}

\subsection*{Your turn}
\begin{itemize}
\raggedright
  \item \emph{Say it:} \emph{uksānā}
  \item \emph{Say it:} \emph{uksānā}, once more
  \item \emph{Say it:} \emph{us ne logoṅ ko zulm ke khilāf uksāyā.} — He stirred people up against oppression.
  \item \emph{Recall:} say the Urdu for favouritism, then the Urdu for oppression, then say \emph{uksānā} again
\end{itemize}

\begin{tcolorbox}[breakable,colback=teal!4,colframe=teal!35!black,title={Before you move on}]
What does \ur{اکسانا} mean? (\textquotedblleft{}To provoke; to stir up\textquotedblright{}.) Say it once more.
\end{tcolorbox}
